\DHead{I. 时代呼唤：迈向智能协作新纪元}
我们正身处人类文明从“计算互联”向“智能共生”跃迁的关键节点。人工智能技术的爆发式增长，特别是大模型与 AI 智能体的崛起，以前所未有的深度重塑着全球经济与社会结构。然而，全球智能鸿沟持续扩大，技术主权博弈加剧，数据孤岛与算法黑盒成为悬在全人类头顶的达摩克利斯之剑。面对气候变化、公共卫生挑战以及复杂的国际关系，唯有坚持开源共享、协作共治，才能确保 AI 向善，避免技术垄断带来的不平等。

在此，我们 —— 来自中国政府、世界顶尖科技企业、国际学术机构、开源基金会及全球开发者社区的代表——郑重发布本宣言。我们呼吁以开放精神为指引，依托世界人工智能合作组织（WAICO）等多边平台，共建共享一个更加公平、透明、可持续的智能未来。

本宣言提出全球 AI 开源行动框架：以数字基础设施为基石，以实质贡献与伦理治理为支柱，以社会价值实现为使命，推动开源软件成为驱动智能时代可持续发展的核心引擎。

\DHead{II. 共同愿景：构建智能共治的数字文明}
我们致力于构建一个开放、包容、安全、可持续的全球 AI 开源生态，使强大的智能技术真正服务于全人类福祉。

我们坚信：开源不仅是代码的共享，更是信任的建立。在生成式 AI 与大模型时代，开源构成了全球数字公共品的底层架构。唯有坚持多元参与、透明可解释、贡献导向的协作原则，才能让技术创新成为推动和平、繁荣与可持续发展的共同力量。

我们共同倡导：将开源创新纳入全球数字公共品战略，作为实现《2030 年可持续发展议程》的关键路径，助力构建共治、共享、共赢的数字文明。
